\documentclass[11pt]{article}

\usepackage[utf8]{inputenc}
\usepackage[T1]{fontenc}
\usepackage{geometry}
\geometry{margin=1in}
\usepackage{microtype}
\usepackage{amsmath,amssymb,amsthm,mathtools}
\usepackage{booktabs}
\usepackage{graphicx}
\usepackage{hyperref}
\usepackage[numbers]{natbib}
\usepackage[capitalize,noabbrev]{cleveref}

\newtheorem{definition}{Definition}[section]
\newtheorem{proposition}{Proposition}[section]
\newtheorem{remark}{Remark}[section]

\newcommand{\buyers}{\mathcal{B}}
\newcommand{\sellers}{\mathcal{S}}
\newcommand{\edges}{\mathcal{E}}
\newcommand{\R}{\mathbb{R}}
\newcommand{\E}{\mathbb{E}}
\newcommand{\inner}[2]{\langle #1, #2 \rangle}

\title{Graph Foundation Models for Two-Sided Market Matching\\
\large Working Paper and Experimental Scaffold}
\author{Working draft}
\date{\today}

\begin{document}
\maketitle

\begin{abstract}
Classical matching theory gives strong guarantees once agents' preference
orderings are known, but real platforms do not observe preferences directly.
They observe timestamped interaction traces on dynamic bipartite graphs. This
paper develops a working bridge between those two views. We formalize the market
as a temporal bipartite multigraph, study preference recovery from typed
interaction histories, and compose recovered scores with both soft and hard
matching operators. The current contribution is intentionally pragmatic: a
complete LaTeX paper draft together with a runnable synthetic experiment stack.
The implementation compares a pointwise gradient-boosted baseline against a
compact graph-matching model that combines recency-weighted typed aggregation, a
Sinkhorn relaxation, and a blocking-pair penalty. On the current synthetic run,
the graph-aware model improves buyer-side Kendall-$\tau$, stability ratio, and
social welfare relative to the pointwise baseline, while still falling short of
the oracle Gale--Shapley upper bound. The paper therefore serves both as a
research statement and as a concrete foundation for a larger temporal graph
foundation model programme.
\end{abstract}

\section{Introduction}

Two-sided matching theory, beginning with Gale and Shapley, studies markets in
which agents on one side are matched to agents on another side under latent
preferences \citep{gale1962college}. In the classical formulation, those
preferences are given. In operational platforms such as food delivery,
prediction markets, labor marketplaces, and acquisition search, they are not.
Instead, the observable object is a history of timestamped interactions:
views, messages, orders, bids, fills, cancels, and repeated engagements.

Modern industrial matching stacks often collapse this history into a pointwise
prediction problem. A tabular model scores buyer--seller pairs independently,
and a matching routine is applied only after those scores have already been
produced. That separation is convenient, but it throws away exactly the
relational structure that matching markets generate: repeated encounters,
typed edges, temporal recency, market thickness, and cross-agent exposure.

This paper works toward a graph-native alternative. The central hypothesis is
that a temporal graph representation can recover latent preference structure
well enough to improve the economic properties of the induced matching, not
just its ranking metrics. That requires more than link prediction. The
evaluation target must include stability and welfare, because the output is a
market allocation rather than a list of pairwise probabilities.

The project in this repository is the first executable milestone for that
larger agenda. It does three things:
\begin{enumerate}
  \item formalizes a two-sided market as a temporal bipartite multigraph;
  \item implements a compact graph-matching model that approximates the paper's
  preference-recovery and matching pipeline;
  \item evaluates the resulting matchings against a pointwise baseline and an
  oracle stable-matching upper bound.
\end{enumerate}

\section{Problem Formulation}

\begin{definition}[Temporal bipartite multigraph]
A temporal bipartite multigraph is a tuple
\[
G = (\buyers, \sellers, \edges, \mathcal{R}, \tau, \phi, \psi),
\]
where $\buyers$ and $\sellers$ are disjoint finite node sets, $\mathcal{R}$ is
a finite set of edge types, $\edges \subseteq \buyers \times \sellers \times
\mathcal{R} \times \R_{\ge 0}$ is the set of typed timestamped edges, $\tau$
extracts edge timestamps, $\phi$ maps nodes to features, and $\psi$ maps edges
to edge attributes.
\end{definition}

The proposal motivating this paper focuses on recovering latent preference
profiles from the observed graph and then feeding those preferences into a
matching operator. Let $u_b(s)$ denote buyer $b$'s latent utility for seller
$s$ and $u_s(b)$ the seller-side analogue. Those utilities induce preference
orderings but are not directly observed.

\begin{definition}[Preference recovery map]
A preference recovery map is a measurable operator
\[
\Pi : G^{(\le t)} \mapsto \succ^{(t)},
\]
from the graph history up to time $t$ to a profile of buyer-side and
seller-side preference orderings.
\end{definition}

This framing highlights the core estimation question: can a temporal graph
model recover the ordinal structure of the latent utilities closely enough that
the resulting matching behaves as if the true preferences had been observed?

\subsection{Matching and Stability}

Given buyer-side scores $\hat U^{B} \in \R^{|\buyers| \times |\sellers|}$ and
seller-side scores $\hat U^{S} \in \R^{|\sellers| \times |\buyers|}$, we study
one-to-one matchings. A hard matching is computed with deferred acceptance,
while a differentiable relaxation is computed with Sinkhorn normalization.

\begin{definition}[Blocking pair]
Let $\mu$ be a matching. A pair $(b,s)$ is blocking under latent utilities if
\[
u_b(s) > u_b(\mu(b))
\quad\text{and}\quad
u_s(b) > u_s(\mu^{-1}(s)),
\]
where unmatched partners are treated as utility $-\infty$.
\end{definition}

\begin{remark}
For the current synthetic scaffold, blocking-pair counts are computed against
the latent utilities used to generate the market. In real data this would need
to be approximated through outcome proxies or structural assumptions.
\end{remark}

\begin{proposition}[Approximate stability via preference error]
Suppose the recovered preference ordering for each buyer and seller differs from
the latent ordering by at most $k$ pairwise inversions, and suppose rank
displacements translate to utility error at Lipschitz rate $L / n$, where $n$
is the market size. Then a matching that is stable under the recovered
preferences is $O(Lk / n)$-approximately stable under the latent utilities.
\end{proposition}

\noindent\textbf{Proof sketch.}
Each inversion perturbs the local rank only finitely, so the utility of the
matched partner can move by at most $O(Lk / n)$ under the assumed rank-to-value
regularity. A blocking pair under the latent utilities must therefore produce
only a bounded improvement unless the recovered preference ordering has made a
much larger error. The proposition is informal, but it captures the reason to
measure Kendall-$\tau$ in the experiment section.

\section{Modeling Approach}

\subsection{Compact Graph-Matching Model}

The long-term target is a temporal graph foundation model. The current
implementation is intentionally smaller. It aggregates typed buyer--seller
interactions with an exponential recency decay, forms buyer and seller context
vectors, and combines those with observed node attributes into low-dimensional
embeddings. The resulting score matrices are
\[
\hat U^{B}_{bs} = f_B(x_b, x_s, h_{bs}),
\qquad
\hat U^{S}_{sb} = f_S(x_s, x_b, h_{bs}),
\]
where $x_b$ and $x_s$ are buyer and seller features and $h_{bs}$ denotes the
typed history tensor for the pair.

This design is not yet a full message-passing network. It is closer to a
hand-parameterized graph layer with shared typed aggregators. That is a feature,
not a bug, for the current stage. It keeps the paper's end-to-end logic
transparent while still testing the actual market-level claim.

\subsection{Soft Matching and Stability Loss}

Given a combined affinity matrix $M$, the implementation computes a soft
matching matrix $P$ with Sinkhorn normalization, following the
entropy-regularized optimal transport perspective \citep{cuturi2013sinkhorn}.
The objective is a sum of:
\begin{enumerate}
  \item a reconstruction term aligning the learned scores with observed
  recency-weighted interactions;
  \item a blocking-pair penalty that suppresses mutually beneficial deviations;
  \item a welfare term that rewards matchings with higher joint utility.
\end{enumerate}

That objective is a tractable proxy for the larger proposal's idea: train the
representation and matching layers jointly rather than treating matching as an
afterthought after pair scoring.

\section{Synthetic Experimental Design}

\subsection{Market Generator}

The synthetic market implemented in this repository generates buyer and seller
feature vectors, latent utility matrices, and typed temporal interactions. The
current edge types are \texttt{view}, \texttt{message}, and \texttt{order}.
Buyers sample sellers according to a softmax over their latent utilities with a
small exposure term, and richer edge types are emitted with higher probability
when the joint buyer-side and seller-side utility is high.

This synthetic setup is deliberately narrow:
\begin{itemize}
  \item one-to-one matching only;
  \item equal market size on both sides;
  \item a single market rather than a pretraining corpus;
  \item evaluation against known latent utilities.
\end{itemize}

The narrowness is useful. It lets us ask the paper's main question before we
invest in the heavier ChronoGraph-style stack: do graph-aware recovered
preferences improve matching outcomes relative to pointwise scoring?

\subsection{Evaluation Protocol}

The current experiment compares three systems:
\begin{enumerate}
  \item \texttt{pointwise\_gbm}: a tabular baseline trained on pair features and
  aggregate history statistics;
  \item \texttt{compact\_graph\_matcher}: the graph-aware scorer with Sinkhorn
  and blocking-pair regularization;
  \item \texttt{oracle\_gale\_shapley}: deferred acceptance on the true latent
  buyer-side and seller-side utilities.
\end{enumerate}

Metrics include:
\begin{itemize}
  \item buyer-side and seller-side Kendall-$\tau$ against the latent utilities;
  \item a true blocking-pair stability ratio;
  \item latent social welfare of the induced hard matching;
  \item MRR and future-MRR ranking diagnostics.
\end{itemize}

\section{Current Results}

\IfFileExists{generated_results.tex}{
  \input{generated_results.tex}
}{
  \paragraph{Missing generated results.}
  Run \texttt{python3 experiments/render\_empirical\_note.py} from the project
  root to materialize the current tables and figure for this section.
}

\section{Discussion}

The current results are encouraging for a first scaffold, but they should be
read carefully. The compact graph-matching model is still much closer to a
structured research baseline than to a modern graph foundation model. It has no
cross-market pretraining, no learned temporal attention, and no explicit
many-to-one capacity constraints. The current experiment is therefore evidence
for the paper's \emph{direction}, not yet for its final strongest claim.

The main conceptual takeaway is nonetheless clear. The pointwise baseline is
competitive on pair prediction, but it does not optimize the market object we
ultimately care about. Once we evaluate at the level of induced matchings,
graph-aware preference recovery already looks more appropriate. The gap between
the compact matcher and the oracle also clarifies the empirical opportunity:
there is enough headroom to justify the more ambitious temporal GFM programme.

\section{Next Steps}

The immediate next steps are straightforward:
\begin{enumerate}
  \item replace the compact aggregator with a true temporal GNN or temporal GFM
  encoder;
  \item add the missing graph baselines, especially static bipartite graph
  models and temporal graph baselines such as TGN-style encoders
  \citep{rossi2020temporal,trivedi2019dyrep};
  \item port the evaluation to a real trader--market slice from the Polymarket
  stack in the broader repository;
  \item extend the matching operator to many-to-one and many-to-many settings,
  where the classical stability notion becomes more subtle
  \citep{rothsotomayor1990two}.
\end{enumerate}

\bibliographystyle{plainnat}
\bibliography{references}

\end{document}
